\documentclass[10pt]{article}
\usepackage{icmlko}

\icmlmeta{IP 파운데이션 월드모델}{166}{같은 축을 세 번 재고 세 번 다른 답을 얻었다}{제작사 이력의 5/5가 우연인지 보려고 세 도메인을 더 붙였다. 5/8이 나왔고, 0/8이 나왔고, 7/8과 8/8이 나왔다 --- 축은 그대로였고 붙이는 방식만 달랐다}

\begin{document}

\twocolumn[
\icmlheader

\begin{abstract}
노트 165가 ``제작사 이력''을 따로 세워 사전 부호 \textbf{5/5} 를 얻고
``다섯 칸뿐이라 우연일 수 있다''를 남겼다. 안 붙은 세 도메인(게임 · 애니 ·
웹툰)에 같은 구성물을 만들어 여덟 칸으로 늘렸다 --- 각 도메인의 작가 ·
제작사 · 퍼블리셔로 \emph{표본 안 시간 인과} 를 세는, 기존 일곱과 같은
방식이다. 덮음 100\%가 셋 다 나왔다. 그리고 \textbf{답이 세 번 달랐다.}
\textbf{① 복제로 붙이니 F18 배깅 5/8, F6 능형 \emph{0/8}} --- 능형은 여덟
도메인 전부 음수이고 크기도 $-$0.014$\sim$$-$0.041로 크다. 사전 부호가
통째로 틀린 것처럼 보인다. \textbf{② 그런데 새로 만든 셋만 X 이고 기존
다섯은 전부 O 였다} --- 출처로 완벽히 갈린다. 그게 단서였다. 기존 다섯은
\texttt{venue\_prominence} 슬롯에 \emph{그대로 남아 있고} 새 축에
복제됐으므로 같은 열이 둘이다. \textbf{③ 옮기기로 바꾸니 F18 7/8 · F6
8/8} 이다. 복제를 걷어내자 능형이 0/8에서 8/8로 뒤집힌다. 원자료도
같은 말을 한다 --- 여덟 도메인 중 일곱에서 라벨과 양의 상관이고(게임
$+$0.301 · 아이돌 $+$0.474 · 모바일 $+$0.212) 펀딩만 $-$0.207이다.
\textbf{노트 165의 5/5 도 복제 아래에서 잰 것이라 폐기한다.} 남는 값은
둘이다 --- 깨끗하게 재면 \textbf{제작사 이력은 여덟 도메인에서 7/8 · 8/8
로 움직이는 축}이고, 그것을 슬롯에서 꺼내는 값은 \textbf{판 $\rho$
$-$0.0104 $\pm$ 0.0032($t{=}3.2$, 갈린다)} 다. 그리고 방법 쪽 교훈 하나 ---
이 흐름에서 측정 방식이 헤드라인을 뒤집은 것이 \textbf{세 번째}다(노트 156
순위 출력, 노트 160 축 방향, 이번 복제).
\end{abstract}
]

\begin{figure*}[t]
\centering
\includegraphics[width=\textwidth]{figs/three.pdf}
\caption{왼쪽 --- 같은 축, 세 가지 답. 가운데 --- 원자료는 여덟 중 일곱이
양수. 오른쪽 --- 옮기는 값은 F18 에서만 갈린다.}
\end{figure*}

\section{세 도메인을 더 붙였다}

\texttt{lab/creatoraxes.py}. 웹툰은 작가, 애니는 제작사, 게임은 퍼블리셔로
\emph{같은 주체 · 더 이른 날짜} 레코드를 센다 --- 기존 일곱 모듈이 쓰는
``표본 안 시간 인과''와 같은 방식이라 사전 정보이고 노트 141의 시점 규칙을
통과한다.

\begin{table}[t]
\centering\tsize
\caption{새로 만든 셋. 덮음 100\%.}
\begin{tabular}{lrr}
\toprule
& 세는 것 & 사전작 0건 \\
\midrule
웹툰 & 작가 & 60.6\% \\
애니 & 제작사 & 39.7\% \\
게임 & 퍼블리셔 & 61.4\% \\
\bottomrule
\end{tabular}
\end{table}

\section{답이 세 번 달랐다}

\begin{table}[t]
\centering\tsize
\caption{같은 축, 붙이는 방식만 다르다.}
\begin{tabular}{llr}
\toprule
& 방식 & 사전 부호 \\
\midrule
노트 165 & 복제 (7도메인) & 5/5 \\
이 노트 & 복제 (8) · F18 & 5/8 \\
이 노트 & 복제 (8) · F6 & \textbf{0/8} \\
\midrule
\textbf{이 노트} & \textbf{옮김 (8) · F18} & \textbf{7/8} \\
\textbf{이 노트} & \textbf{옮김 (8) · F6} & \textbf{8/8} \\
\bottomrule
\end{tabular}
\end{table}

\claimx{\textbf{복제에서는 능형이 0/8이다.} 여덟 도메인 전부 음수이고
크기가 $-$0.014$\sim$$-$0.041로 작지도 않다. 이것만 보면 ``사전 부호가
통째로 틀렸다''로 읽힌다.}{실제로 처음에 그렇게 읽을 뻔했다. 노트 160에서
입장 마찰로 똑같은 함정에 빠졌던 터라 한 번 더 갈랐다.}

\claimx{\textbf{단서는 출처로 완벽히 갈린다는 것이었다.} F18 복제에서 새로
만든 셋(게임 · 애니 · 웹툰)이 전부 X 이고 기존 다섯이 전부 O 다. 축이
문제라면 출처와 무관해야 한다.}{기존 다섯은 \texttt{venue\_prominence}
슬롯에 남아 있으면서 새 축에 복제됐다 --- \textbf{같은 열이 둘}이다.
노트 165가 잰 공선성이 여기서 부호까지 뒤집었다.}

\claimx{\textbf{옮기기로 바꾸니 F18 7/8 · F6 8/8 이다.} 복제를 걷어내자
능형이 0/8에서 8/8로 뒤집힌다. 유일한 X 는 F18 의 게임($-$0.0070)이다.}{원자료도
같은 말을 한다 --- 여덟 중 일곱이 양의 상관이고 펀딩만 $-$0.207 이다.
\textbf{사전 부호 $+1$ 이 맞다.}}

\begin{table}[t]
\centering\tsize
\caption{creator\_track 과 라벨의 상관. 유보.}
\begin{tabular}{lrl}
\toprule
아이돌 & $+$.474 & \\
\textbf{게임} & \textbf{$+$.301} & 새로 만든 것 \\
모바일 & $+$.212 & \\
\textbf{웹툰} & \textbf{$+$.105} & 새로 만든 것 \\
\textbf{애니} & \textbf{$+$.072} & 새로 만든 것 \\
도서 & $+$.062 & \\
세계애니 & $+$.056 & \\
\textbf{펀딩} & \textbf{$-$.207} & 유일한 음수 \\
\bottomrule
\end{tabular}
\end{table}

\section{그래서 남는 값 둘}

\claimx{\textbf{① 제작사 이력은 여덟 도메인에서 7/8 · 8/8 로 움직인다.}
노트 163의 타깃 폭(45/45) · 굿즈 규모(39/40) 다음이고, 셋 다 구성물이
하나인 축이다.}{노트 165의 5/5 는 폐기한다 --- 복제 아래에서 잰 값이다.
\textbf{같은 축이 5/5 에서 0/8 을 거쳐 8/8 로 갔다.}}

\claimx{\textbf{② 슬롯에서 꺼내는 값은 판 $\rho$ $-$0.0104 $\pm$
0.0032($t{=}3.2$)다} --- F18 에서만 갈리고 F6 은 $-$0.0012($t{=}0.5$)로
안 갈린다.}{노트 165가 일곱 도메인으로 잰 $-$0.0122 와 거의 같다. 세
도메인을 더 붙여도 비용은 안 줄었다 --- \textbf{비용은 새 축이 얇아서가
아니라 옛 슬롯이 빌려오기를 잃어서다.}}

\section{방법 쪽 교훈}

\claimx{\textbf{이 흐름에서 측정 방식이 헤드라인을 뒤집은 것이 세
번째다.} 노트 156(순위를 내는 정식화에 평균 차를 물었다) · 노트 160(라벨로
방향을 정한 열에 원 방향 사전 부호를 댔다) · 이번(복제로 공선성을
만들었다).}{셋 다 ``모형이 이상하다''로 읽힐 뻔했고 셋 다 재는 쪽 문제였다.
\textbf{모형이 상식과 크게 어긋나면 먼저 재는 자를 의심한다} --- 이 판에서
세 번 연속 그랬다.}

\begin{table}[t]
\centering\tsize
\caption{남은 것.}
\begin{tabular}{ll}
\toprule
펀딩 & 유일한 음수($-$.207) --- 창작자가 많을수록 나쁜가 \\
채택 & 판 $-$.0104 대 축 하나 --- 무엇과 바꾸나 \\
\textbf{규약} & 축을 더할 때 원래 자리에 남기지 말 것 \\
\bottomrule
\end{tabular}
\end{table}

마지막 줄을 규약으로 적어 둔다. \textbf{축을 새로 세울 때 원래 슬롯에
같이 남기면 공선성이 해석을 뒤집는다} --- 예측은 거의 안 변하는데
(노트 165: 판 $-$0.0029, $t{=}1.6$) 부호는 8/8이 0/8이 된다.

\end{document}
